\section{Introduction}
\label{sec:introduction}

\subsection{Motivation}
\label{subsec:motivation}
The final approach and landing is statistically the most accident-prone phase of
flight. Guiding it today relies on ground infrastructure such as instrument landing
systems and on satellite navigation, which are not available at every runway and can
be degraded or spoofed. A forward-looking camera on the aircraft offers a low-cost,
self-contained alternative: if the runway can be located in the image and its
geometry recovered, the system can estimate the aircraft's relative pose and support
alignment during approach, reducing pilot workload and enabling assisted or
autonomous landing~\cite{ferreira2025landing}.

Detecting the runway from the aircraft perspective is, however, a hard vision
problem. The runway appears as a small, thin quadrilateral early in the approach and
grows and shifts rapidly as the aircraft descends, so its \emph{apparent size and
distance} vary over orders of magnitude. Its \emph{orientation and perspective}
depend on the approach angle, glide slope and lateral offset, and the on-board
\emph{camera viewpoint} adds further variability. On top of this geometric
variability, \emph{weather and illumination} --- rain, sun glare, fog, snow and night
operations --- degrade contrast and inject nuisances (streaks, halos, scattering,
low light) precisely when reliable detection matters most. Runway-specific detectors
such as YOLO-RWY~\cite{yolorwy2024} and lightweight tracking-based landing-guidance
pipelines~\cite{ferreira2025landing, runwaysurveillance2026} illustrate both the
practical interest in the task and its difficulty.

Framing the task as \emph{keypoint (pose) estimation} of the four runway corners,
rather than as plain bounding-box detection, is particularly attractive: the four
corners directly parameterize the runway quadrilateral and can be fed into downstream
pose-from-geometry estimators, while remaining cheap enough for real-time, on-board
inference.

The central obstacle is training data. Annotated real approach imagery that spans a
wide range of airports, runway geometries, approach angles and, crucially, weather
and lighting conditions is scarce: real landing footage is costly to collect and its
per-corner annotation is laborious, so publicly available labeled datasets are
limited. Photorealistic simulators and geospatial rendering pipelines offer a way
out --- they synthesize large, automatically and exactly labeled datasets covering
many locations and viewpoints~\cite{syntheticrunway2025} --- and the Landing Approach
Runway Detection (LARD) dataset family~\cite{lard2023} exemplifies this approach. Its
second iteration (LARD-V2) provides several independent synthetic image sources for
the same set of runway locations.

\subsection{Problem Statement}
\label{subsec:problem}
Models trained purely on synthetic imagery rarely reach their synthetic-domain
performance once deployed on real images. This performance drop, commonly referred
to as the \emph{Sim2Real gap}, is a specific instance of \emph{domain shift}:
the statistical distribution of the training (source, synthetic) domain differs
from that of the deployment (target, real) domain in texture, lighting, sensor
noise, atmospheric effects and scene composition. For runway detection the gap is
amplified by safety-critical operating conditions: the system must remain reliable
precisely in the adverse weather and low-visibility scenarios (rain, glare, fog,
snow, night) that are hardest to render faithfully and least represented in nominal
synthetic data.

This thesis studies the Sim2Real gap for runway keypoint detection when a
YOLOv8-pose model is fine-tuned on nominal-condition synthetic data from LARD-V2
and evaluated on a real, weather-annotated test set from LARD. It first
\emph{quantifies} the gap and then investigates three complementary families of
techniques for \emph{closing} it, ordered by how directly they intervene on the
domain shift: input-space weather augmentation, self-training with pseudo-labels
on unlabeled real footage, and generative style transfer that imports
target-domain appearance onto the synthetic source.

\paragraph{Goal of the thesis.}
The goal of this thesis is to evaluate \emph{which Sim2Real strategy best improves
lightweight runway-corner detection from aircraft-view images when only limited
annotated real data is available}. Concretely, the strategies are compared under a
single, reproducible, lightweight setup (a fine-tuned YOLOv8-pose model, a fixed
5000-image training budget, and a shared weather-stratified real test set), so that
observed differences are attributable to the adaptation strategy rather than to
model capacity or training budget.

\paragraph{Research gap.}
Prior work tends to address these ingredients \emph{separately}: some studies focus
on generating and mixing synthetic runway data~\cite{syntheticrunway2025}, others on
runway-specific detector design~\cite{yolorwy2024}, weather-robust
augmentation~\cite{weatherrobust2025}, or domain adaptation~\cite{ganin2015dann}. In
contrast, this thesis compares several Sim2Real strategies --- weather augmentation,
self-training, and diffusion-based style transfer --- head-to-head under one
lightweight, reproducible protocol, which makes their relative benefits directly
comparable.

\paragraph{Label-preserving style transfer.}
A recurring difficulty for the generative strategy is that the four runway corners
must remain geometrically valid after an image is restyled. Unpaired image-to-image
translation with CycleGAN-style models can hallucinate or displace structure, which
would silently corrupt the corner labels. This thesis therefore favors a
\emph{conditional diffusion} model that is guided by the source image's structure so
that it repaints appearance --- lighting, texture, weather traits of the real target
domain --- while leaving the runway geometry, and hence the labels, in place; the
CycleGAN route is kept only as a contrast that exposes the label-displacement risk.

\subsection{Research Questions}
\label{subsec:research-questions}
The work is organized around four research questions (RQs) that follow a single
logical thread: measure the gap, then progressively attack it. Each RQ states a
\emph{question}; the \emph{experiment} that answers it is developed in the
corresponding section (\Cref{sec:gap,sec:weather-aug,sec:pseudo-labels,sec:style-transfer}),
keeping the questions separate from the methodology.
\begin{itemize}
    \item \textbf{RQ1 (Gap characterization).} How large is the Sim2Real gap for
    runway keypoint detection across the five LARD-V2 synthetic sources and a
    combined source, and which source generalizes best to the real test set
    overall and within each individual weather scenario? \emph{Motivation:} the size
    and structure of the gap must be measured before any mitigation can be assessed,
    and the choice of synthetic source is itself a design decision reported to matter
    for runway detection~\cite{syntheticrunway2025}.
    \item \textbf{RQ2 (Weather augmentation).} To what extent do input-space
    weather augmentations improve real-world generalization for a single synthetic
    source and for the combined source? Which augmentation ratios help and which
    harm, and can augmentation compensate for a lack of dataset diversity?
    \emph{Motivation:} augmentation is the cheapest, label-preserving robustness
    lever and is known to shift detector performance under adverse
    weather~\cite{weatherrobust2025}, but its ratio and its interaction with dataset
    diversity are rarely quantified for this task.
    \item \textbf{RQ3 (Self-training).} Can self-training with confidence-filtered
    pseudo-labels on unlabeled real landing sequences reduce the gap, and does
    exploiting temporal consistency through an object tracker further improve the
    pseudo-labels? \emph{Motivation:} unlabeled real landing footage is comparatively
    abundant, and self-training is a standard route to unsupervised domain
    adaptation~\cite{ganin2015dann}; tracking is already used in landing-guidance
    pipelines~\cite{ferreira2025landing, runwaysurveillance2026} and can enforce
    temporal label consistency.
    \item \textbf{RQ4 (Style transfer).} Can diffusion-based style transfer import
    real-domain (and adverse-weather) appearance onto the synthetic source without
    breaking the geometric labels, and how does it compare to a GAN-based
    alternative that risks displacing keypoints? \emph{Motivation:} style transfer
    attacks the appearance gap most directly, but only helps if the four-corner
    labels survive the transformation.
\end{itemize}

\subsection{Contributions}
\label{subsec:contributions}
This thesis makes the following contributions:
\begin{enumerate}
    \item A \textbf{real-domain characterization of the Sim2Real gap} for runway
    keypoint detection, using a shared-location cross-evaluation across five
    synthetic rendering sources plus a combined source against the real LARD-V1 set.
    This complements the LARD-V2 benchmark (which studied cross-dataset
    generalization on \emph{synthetic} test sets) and isolates the rendering source
    from the runway viewpoint.
    \item A \textbf{vertical-angle test-set filtering} method and the finding that
    the gap is only measurable once the test sets are matched on runway
    \emph{orientation} as well as size; without it, the real set spuriously appears
    easier than the synthetic sets. The residual synthetic difficulty is shown to be
    largely an image-quality effect concentrated on distant runways.
    \item A \textbf{weather-stratified real benchmark}, extended with private landing
    sequences to compensate for the sparse adverse-weather representation, together
    with a style-feature analysis explaining why different sources suit different
    scenarios.
    \item A systematic \textbf{study of weather augmentations} (glare, fog, rain,
    water refraction, night) and their mixing ratios, with a variance-controlled
    protocol, testing whether augmentation substitutes for dataset diversity.
    \item A \textbf{self-training pipeline} on unlabeled real landing sequences that
    compares confidence thresholds, analyzes pseudo-label composition, and extends
    pseudo-labeling with a tracker-based temporal-consistency mechanism (ongoing).
    \item An investigation of \textbf{control-guided diffusion style transfer}
    (Cosmos Transfer) for label-preserving domain translation of the synthetic
    source, with an optional CycleGAN comparison (ongoing).
\end{enumerate}

\subsection{Thesis Outline}
\label{subsec:outline}
\Cref{sec:background} introduces the concepts required to follow the work: vision-based
runway detection, YOLO pose estimation, the Sim2Real gap and domain-adaptation
families, and the LARD datasets. \Cref{sec:related-work} positions the thesis
against related literature. \Cref{sec:setup} defines the shared experimental
setup: datasets and weather splits, model, evaluation metrics and training
protocol. \Cref{sec:gap} establishes and dissects the Sim2Real gap (RQ1).
\Cref{sec:weather-aug} studies weather augmentations (RQ2). \Cref{sec:pseudo-labels}
develops the self-training pipelines (RQ3). \Cref{sec:style-transfer} investigates
diffusion-based style transfer (RQ4). \Cref{sec:discussion} compares the techniques
and discusses limitations, and \Cref{sec:conclusion} concludes and outlines future
work.
